\section{Experiments}
\label{sec:experiments}

Open with the questions your experiments answer, for example:
(1) Does \method outperform prior methods? (2) Does it generalize to unseen objects and scenes? (3) Which components matter most?
Then answer each question in its own subsection.

\subsection{Setup}
Describe the benchmarks, robot platforms, baselines, metrics, and the number of trials per task, so that results can be reproduced. Put full details (hyperparameters, task lists) in the appendix.

\subsection{Main Results}

% Sample results table kept in its own file under tables/ and pulled into the
% evaluation section with % Sample results table kept in its own file under tables/ and pulled into the
% evaluation section with % Sample results table kept in its own file under tables/ and pulled into the
% evaluation section with \input{tables/main_results}. Keeping tables in a
% separate folder keeps your section prose readable.
\begin{table}[t]
    \centering
    \begin{tabular}{lcc}
        \toprule
        Method & Metric A $\uparrow$ & Metric B $\uparrow$ \\
        \midrule
        Baseline A           & 70.0          & 61.2 \\
        Baseline B           & 74.3          & 68.0 \\
        \methodname{} (ours) & \textbf{82.5} & \textbf{74.8} \\
        \bottomrule
    \end{tabular}
    \caption{\textbf{Main results.} Bold marks the best result per column.
    Replace with your own numbers and add standard errors where appropriate.}
    \label{tab:main-results}
\end{table}
. Keeping tables in a
% separate folder keeps your section prose readable.
\begin{table}[t]
    \centering
    \begin{tabular}{lcc}
        \toprule
        Method & Metric A $\uparrow$ & Metric B $\uparrow$ \\
        \midrule
        Baseline A           & 70.0          & 61.2 \\
        Baseline B           & 74.3          & 68.0 \\
        \methodname{} (ours) & \textbf{82.5} & \textbf{74.8} \\
        \bottomrule
    \end{tabular}
    \caption{\textbf{Main results.} Bold marks the best result per column.
    Replace with your own numbers and add standard errors where appropriate.}
    \label{tab:main-results}
\end{table}
. Keeping tables in a
% separate folder keeps your section prose readable.
\begin{table}[t]
    \centering
    \begin{tabular}{lcc}
        \toprule
        Method & Metric A $\uparrow$ & Metric B $\uparrow$ \\
        \midrule
        Baseline A           & 70.0          & 61.2 \\
        Baseline B           & 74.3          & 68.0 \\
        \methodname{} (ours) & \textbf{82.5} & \textbf{74.8} \\
        \bottomrule
    \end{tabular}
    \caption{\textbf{Main results.} Bold marks the best result per column.
    Replace with your own numbers and add standard errors where appropriate.}
    \label{tab:main-results}
\end{table}


\Cref{tab:main} shows an example table. Use \texttt{booktabs} rules (\verb|\toprule|, \verb|\midrule|, \verb|\bottomrule|) and no vertical lines, put the caption above the table, bold the best result, and shade your own row (see \texttt{tables/main\_results.tex}). Report the number of trials or a confidence interval with every success rate.

\begin{figure}[t]
    \centering
    \begin{subfigure}[b]{0.48\linewidth}
        \centering
        \includegraphics[width=\linewidth]{example-image-b}
        \caption{First panel.}
        \label{fig:panel_a}
    \end{subfigure}
    \hfill
    \begin{subfigure}[b]{0.48\linewidth}
        \centering
        \includegraphics[width=\linewidth]{example-image-c}
        \caption{Second panel.}
        \label{fig:panel_b}
    \end{subfigure}
    \caption{\textbf{A figure with sub-figures.} Use the \texttt{subcaption} package for multi-panel figures, and save plots as vector PDFs so that they stay sharp.}
    \label{fig:panels}
\end{figure}

\subsection{Ablations}
Remove or replace one component of \method at a time and report the effect, as in \cref{fig:panels}. Explain \emph{why} each result happens, not just what the numbers are.
